\documentclass[11pt]{article}
\usepackage{amsmath,amssymb,amsthm,mathtools,bm}
\usepackage[margin=1in]{geometry}
\usepackage{hyperref}
\usepackage{enumitem}
\title{\textbf{Problem 2 Solution: Local Stability, Linearization, and Feedback Design}}
\author{}
\date{}

\newtheorem{theorem}{Theorem}
\newtheorem{proposition}{Proposition}
\newtheorem{corollary}{Corollary}
\newtheorem{lemma}{Lemma}
\theoremstyle{definition}
\newtheorem{definition}{Definition}
\newtheorem{remark}{Remark}
\newtheorem{assumption}{Assumption}

\begin{document}
\maketitle

\section{Setup}

We work under the translation-only model fixed in Problem 1. The admissible center-of-mass region is
\[
\Omega_r=[r,L-r]^3 \subset \mathbb{R}^3,
\]
and gravity acts as $-mg\,e_z$, where $e_z=(0,0,1)^\top$ points upward.

The translational dynamics are
\[
m\ddot x = F_{\mathrm{mag}}(x,u)-mg\,e_z,
\qquad x\in\Omega_r,\quad u\in U\subset\mathbb{R}^6.
\]
For the main analysis we adopt the control-affine surrogate
\[
F_{\mathrm{mag}}(x,u)=G(x)u=\sum_{i=1}^6 u_i g_i(x),
\]
where $g_i\in C^1(\Omega_r;\mathbb{R}^3)$ and
\[
G(x)=
\begin{bmatrix}
| & | & & |\\
g_1(x) & g_2(x) & \cdots & g_6(x)\\
| & | & & |
\end{bmatrix}
\in \mathbb{R}^{3\times 6}.
\]
We assume there is no passive damping or drag unless explicitly introduced.

Let $x^\star\in\operatorname{int}(\Omega_r)$ be an equilibrium point from Problem 1, so there exists $u^\star\in U$ such that
\[
G(x^\star)u^\star=mg\,e_z.
\]

\section{Linearization About an Equilibrium}

Set
\[
x=x^\star+\delta x,
\qquad
u=u^\star+\delta u,
\qquad
\delta v=\delta \dot x.
\]
We derive the first-order perturbation model.

\begin{theorem}[Linearization about an equilibrium]
Let $x^\star\in\operatorname{int}(\Omega_r)$ and $u^\star\in U$ satisfy
\[
G(x^\star)u^\star=mg\,e_z.
\]
Then the first-order perturbation dynamics are
\[
m\,\delta \ddot x = A_x\,\delta x + B_x\,\delta u,
\]
where
\[
A_x := \sum_{i=1}^6 u_i^\star\,Dg_i(x^\star) \in \mathbb{R}^{3\times 3},
\qquad
B_x := G(x^\star) \in \mathbb{R}^{3\times 6}.
\]
Equivalently, in first-order state form with
\[
\xi=
\begin{bmatrix}
\delta x\\
\delta v
\end{bmatrix},
\]
we have
\[
\dot \xi = A\xi+B\delta u,
\]
with
\[
A=
\begin{bmatrix}
0&I\\
\frac{1}{m}A_x&0
\end{bmatrix},
\qquad
B=
\begin{bmatrix}
0\\
\frac{1}{m}B_x
\end{bmatrix}.
\]
\end{theorem}

\begin{proof}
Expand the dynamics about $(x^\star,u^\star)$:
\[
m(\ddot x^\star+\delta \ddot x)=G(x^\star+\delta x)(u^\star+\delta u)-mg\,e_z.
\]
Since $x^\star$ is an equilibrium, $\ddot x^\star=0$ and
\[
G(x^\star)u^\star=mg\,e_z.
\]
Hence
\[
m\,\delta \ddot x = G(x^\star+\delta x)u^\star + G(x^\star+\delta x)\delta u - G(x^\star)u^\star.
\]
Using the first-order expansion
\[
G(x^\star+\delta x)=G(x^\star)+\sum_{j=1}^3 \frac{\partial G}{\partial x_j}(x^\star)\,\delta x_j + o(\|\delta x\|),
\]
we obtain
\[
G(x^\star+\delta x)u^\star
=
G(x^\star)u^\star
+
\left(\sum_{j=1}^3 \frac{\partial G}{\partial x_j}(x^\star)u^\star\,\delta x_j\right)
+o(\|\delta x\|).
\]
Also,
\[
G(x^\star+\delta x)\delta u = G(x^\star)\delta u + o(\|\delta x\|\,\|\delta u\|).
\]
Discarding higher-order terms and canceling $G(x^\star)u^\star=mg\,e_z$ yields
\[
m\,\delta \ddot x =
\left(\sum_{j=1}^3 \frac{\partial G}{\partial x_j}(x^\star)u^\star\,\delta x_j\right)
+
G(x^\star)\delta u.
\]
Because the $i$th column of $G(x)$ is $g_i(x)$, the Jacobian term may be written as
\[
\sum_{i=1}^6 u_i^\star Dg_i(x^\star)\,\delta x.
\]
Thus
\[
A_x=\sum_{i=1}^6 u_i^\star Dg_i(x^\star),
\qquad
B_x=G(x^\star),
\]
as claimed. The first-order state-space form follows by setting $\delta v=\delta \dot x$.
\end{proof}

\begin{remark}[Interpretation]
The matrix $A_x$ captures how the equilibrium field gradients feed position perturbations back into force, while $B_x$ is the first-order actuator allocation map from input perturbations to net force perturbations.
\end{remark}

\section{Stabilizability Conditions}

Problem 2 asks when the linearized system is locally stabilizable.

\begin{proposition}[Necessary condition via PBH geometry]
Let $(A,B)$ be the linearized pair above. If $\operatorname{rank}(B_x)<3$, then there exists a nonzero vector $n\in\mathbb{R}^3$ such that
\[
n^\top B_x=0.
\]
Consequently, perturbations along the translational direction $n$ receive no first-order control input. Therefore, if $A$ has an unstable mode whose translational component lies in such an uncontrollable direction, then $(A,B)$ is not stabilizable.
\end{proposition}

\begin{proof}
If $\operatorname{rank}(B_x)<3$, then the left nullspace of $B_x$ is nontrivial, so there exists $n\neq 0$ with $n^\top B_x=0$. Since
\[
B=
\begin{bmatrix}
0\\
\frac{1}{m}B_x
\end{bmatrix},
\]
this means that no first-order input can directly generate acceleration in the direction $n$. If $A$ has an unstable eigenvector whose translational component lies in this uncontrollable subspace, then the Popov--Belevitch--Hautus criterion fails, so $(A,B)$ is not stabilizable.
\end{proof}

\begin{remark}
Thus $\operatorname{rank}(B_x)=3$ is not by itself sufficient for stabilizability, but it is the natural full translational authority condition. The exact criterion is the PBH stabilizability test.
\end{remark}

\begin{corollary}[Full translational input authority]
If $\operatorname{rank}(B_x)=3$, then the lower block of $B$ spans all translational acceleration directions.
\end{corollary}

\section{Two-Layer Feedback Architecture}

A key point in this problem is that PID should not be interpreted as six unrelated SISO controllers on the six actuator channels. The correct local architecture is:

\begin{enumerate}[label=\textbf{Step \arabic*:},leftmargin=1.6cm]
\item An outer-loop controller computes a desired net force perturbation $f_{\mathrm{cmd}}\in\mathbb{R}^3$ from $(\delta x,\delta v)$.
\item An allocation map computes an actuator perturbation $\delta u\in\mathbb{R}^6$ satisfying, approximately or exactly,
\[
B_x\delta u=f_{\mathrm{cmd}}.
\]
\end{enumerate}

A standard local PID-like outer loop is
\[
f_{\mathrm{cmd}} = -K_p\,\delta x - K_d\,\delta v - K_i\eta,
\qquad
\dot\eta=\delta x,
\]
where $K_p,K_d,K_i\in\mathbb{R}^{3\times 3}$ are gain matrices.

If one ignores the integral term initially, the simplest allocation is the minimum-norm pseudoinverse solution
\[
\delta u = B_x^\dagger f_{\mathrm{cmd}}.
\]

\begin{remark}[Allocation and saturation]
The pseudoinverse gives the minimum-norm unconstrained actuator perturbation, but it may violate actuator bounds even when the desired force is feasible. In practice one replaces it with a bounded least-squares or quadratic-program allocation step.
\end{remark}

\section{Local Asymptotic Stabilization by PD + Allocation}

We now show that the two-layer architecture yields a stabilizing controller when local translational authority is full.

\begin{theorem}[Local asymptotic stabilization by PD force feedback]
Assume $\operatorname{rank}(B_x)=3$. Let $K_p,K_d\in\mathbb{R}^{3\times 3}$ be symmetric positive definite. Suppose the allocation law enforces
\[
B_x\delta u = -A_x\delta x - K_p\delta x - K_d\delta v.
\]
Then the closed-loop linearized perturbation dynamics reduce to
\[
m\,\delta \ddot x + K_d\delta \dot x + K_p\delta x = 0.
\]
Moreover, the origin $(\delta x,\delta v)=(0,0)$ is locally asymptotically stable.
\end{theorem}

\begin{proof}
Because $\operatorname{rank}(B_x)=3$, the linear map $\delta u\mapsto B_x\delta u$ is surjective onto $\mathbb{R}^3$, so the stated allocation equation is solvable for every $(\delta x,\delta v)$. Substituting it into the linearized dynamics
\[
m\,\delta \ddot x = A_x\delta x + B_x\delta u
\]
yields
\[
m\,\delta \ddot x = A_x\delta x - A_x\delta x - K_p\delta x - K_d\delta v,
\]
so
\[
m\,\delta \ddot x + K_d\delta \dot x + K_p\delta x = 0.
\]
Consider the Lyapunov function
\[
V(\delta x,\delta v)=\frac12 m\,\delta v^\top\delta v + \frac12\delta x^\top K_p\delta x.
\]
Since $m>0$ and $K_p\succ 0$, $V$ is positive definite. Along closed-loop trajectories,
\[
\dot V = m\,\delta v^\top\delta \ddot x + \delta x^\top K_p\delta v.
\]
Using the closed-loop dynamics,
\[
m\,\delta \ddot x = -K_d\delta v - K_p\delta x,
\]
we obtain
\[
\dot V = \delta v^\top(-K_d\delta v - K_p\delta x) + \delta x^\top K_p\delta v = -\delta v^\top K_d\delta v \le 0.
\]
Thus $\dot V$ is negative semidefinite. By LaSalle's invariance principle, trajectories converge to the largest invariant subset of
\[
\{(\delta x,\delta v): \dot V=0\} = \{(\delta x,\delta v): \delta v=0\}.
\]
On this set, $\delta v=0$ for all time. Invariance then requires $\delta \dot v=\delta \ddot x=0$, so the closed-loop dynamics imply
\[
K_p\delta x=0.
\]
Since $K_p\succ 0$, this gives $\delta x=0$. Hence the largest invariant subset of $\{\dot V=0\}$ is the singleton $\{(0,0)\}$. By LaSalle's principle, $(\delta x,\delta v)\to (0,0)$, so the origin is asymptotically stable.
\end{proof}

\begin{remark}
This theorem proves local asymptotic stabilization of the linearized closed-loop system. By Lyapunov's indirect method, asymptotic stability of the linearized closed-loop system implies local asymptotic stability of the corresponding equilibrium of the full nonlinear system.
\end{remark}

\section{PID Interpretation and Integral Action}

The assignment explicitly asks whether PID is sufficient locally. The answer is yes, with the correct MIMO interpretation.

\begin{proposition}[Correct local PID interpretation]
A locally meaningful PID controller for this problem is obtained by choosing a desired net force
\[
f_{\mathrm{cmd}} = -K_p\delta x - K_d\delta v - K_i\eta,
\qquad
\dot\eta=\delta x,
\]
and then solving an allocation problem of the form
\[
B_x\delta u \approx f_{\mathrm{cmd}} - A_x\delta x.
\]
This is the proper multivariable generalization of PID for the six-actuator magnetic levitation problem.
\end{proposition}

\begin{proof}
The linearized translational dynamics are
\[
m\,\delta \ddot x = A_x\delta x + B_x\delta u.
\]
If we desire a target closed-loop perturbation force law
\[
m\,\delta \ddot x = -K_p\delta x - K_d\delta v - K_i\eta,
\]
then $\delta u$ must be chosen so that
\[
B_x\delta u = -A_x\delta x - K_p\delta x - K_d\delta v - K_i\eta.
\]
Thus the control task naturally decomposes into an outer-loop PID force design and an inner-loop actuator allocation. This is exactly the multivariable PID interpretation required by the problem.
\end{proof}

\begin{remark}
Integral action is useful for canceling constant model mismatch and static disturbances, but it must be implemented with anti-windup or constrained allocation when actuator limits are present.
\end{remark}

\section{Comparison with State Feedback}

The PD construction above is not the only local controller. Since the linearized system is
\[
\dot\xi=A\xi+B\delta u,
\]
any stabilizing full-state feedback law
\[
\delta u = -K\xi
\]
may be used whenever $(A,B)$ is stabilizable. In particular, LQR provides a systematic way to choose $K$ and often gives better transient shaping than hand-tuned PID gains. The advantage of the two-layer PD/PID approach is physical interpretability: the outer loop produces a desired force, and the inner loop translates that force into actuator commands.

\section{Extension to Physical Magnetic Models}

Problem 2 also asks how the linearization changes for the more physical force laws.

\subsection*{Permanent or controlled dipole model}
If
\[
B(x,u)=\sum_{i=1}^6 u_i B_i(x),
\qquad
F_{\mathrm{mag}}(x,u)=\nabla(\mu\cdot B(x,u)),
\]
and $\mu$ is fixed, then $F_{\mathrm{mag}}$ remains affine in $u$. The matrices $A_x$ and $B_x$ are obtained exactly as above after replacing $g_i(x)$ by the corresponding force basis fields.

\subsection*{Induced-dipole model}
If
\[
F_{\mathrm{mag}}(x,u)=\frac{\alpha}{2}\nabla\|B(x,u)\|^2,
\]
then the force law is quadratic in $u$. In that case one cannot read off $B_x$ from a control-affine matrix $G(x^\star)$, because no such global $G$ exists. Instead,
\[
A_x = \left.\frac{\partial F_{\mathrm{mag}}}{\partial x}\right|_{(x^\star,u^\star)},
\qquad
B_x = \left.\frac{\partial F_{\mathrm{mag}}}{\partial u}\right|_{(x^\star,u^\star)}
\]
must be computed directly from the Jacobians of the nonlinear force law at equilibrium. Once these Jacobians are computed, the local linearized theory above applies unchanged.

\begin{remark}
Thus the local stabilization theory of Problem 2 is robust across both physical models, provided one linearizes the actual force law at the operating point.
\end{remark}

\section{Final Answer to Problem 2}

The formal solution to Problem 2 is:

\begin{enumerate}[label=\textbf{(\arabic*)},leftmargin=1.6cm]
\item The linearization about an equilibrium $(x^\star,u^\star)$ is
\[
m\,\delta \ddot x = A_x\delta x + B_x\delta u,
\qquad
A_x=\sum_{i=1}^6 u_i^\star Dg_i(x^\star),
\qquad
B_x=G(x^\star).
\]

\item In state form,
\[
\dot\xi=A\xi+B\delta u,
\qquad
A=
\begin{bmatrix}
0&I\\
\frac{1}{m}A_x&0
\end{bmatrix},
\qquad
B=
\begin{bmatrix}
0\\
\frac{1}{m}B_x
\end{bmatrix}.
\]

\item If $\operatorname{rank}(B_x)<3$, then some translational directions receive no first-order input authority. Stabilizability then depends on whether $A$ has unstable modes in those uncontrollable directions, as characterized by the PBH test.

\item If $\operatorname{rank}(B_x)=3$, then there is full local translational force authority, and one can allocate actuator perturbations to realize arbitrary commanded net force perturbations.

\item The correct PID interpretation is two-layer: an outer-loop PID computes a desired force, and an inner-loop allocator computes actuator commands satisfying
\[
B_x\delta u \approx f_{\mathrm{cmd}}-A_x\delta x.
\]

\item With full local translational authority, the allocation law
\[
B_x\delta u=-A_x\delta x-K_p\delta x-K_d\delta v
\]
reduces the closed-loop dynamics to a damped second-order system
\[
m\,\delta \ddot x + K_d\delta \dot x + K_p\delta x = 0,
\]
which is asymptotically stable for $K_p,K_d\succ 0$.

\item In the physically grounded magnetic models, the same local theory applies after computing the Jacobians
\[
A_x=\left.\frac{\partial F_{\mathrm{mag}}}{\partial x}\right|_{(x^\star,u^\star)},
\qquad
B_x=\left.\frac{\partial F_{\mathrm{mag}}}{\partial u}\right|_{(x^\star,u^\star)}.
\]
\end{enumerate}

\end{document}
